\subsection*{F 长期状态转移递推}

代码~\ref{lst:state-transition-probability} 展示长期预测中状态转移概率的递推方式。基础矩阵先给出缓和、维持、升级三类状态之间的常规切换概率，再映射到 logit 空间，由物理压力和 GPR/OVX 综合风险压力对升级倾向进行温和修正；后期缓和不再按日历时间机械触发，而由恢复证据门控。

\begin{paperlisting}{lst:state-transition-probability}{python}{长期状态转移概率递推}
def transition_probabilities(current_state, day_index, risk_constraints,
                             physical_pressure, recovery_evidence):
    base_probs = TRANSITION_MATRIX[STATE_INDEX[current_state]].copy()
    logits = np.log(np.clip(base_probs, 1e-6, 1.0))

    # 市场风险主要影响外推早中期；物理压力来自剩余缺口、SPR、库存和绕道
    early_risk_decay = np.clip(1.0 - max(day_index - 60, 0) / 150.0,
                               0.25, 1.0)
    risk_shift = risk_constraints.combined_pressure * early_risk_decay * 0.75
    physical_shift = np.clip(physical_pressure, 0.0, 1.0) * 0.80
    logits[STATE_INDEX["escalation"]] += 0.16 * risk_shift + 0.24 * physical_shift
    logits[STATE_INDEX["easing"]] -= 0.11 * risk_shift + 0.18 * physical_shift

    if current_state == "escalation":
        logits[STATE_INDEX["escalation"]] += 0.07 * risk_shift + 0.11 * physical_shift
    elif current_state == "easing":
        logits[STATE_INDEX["neutral"]] += 0.04 * risk_shift

    # 后期缓和由恢复证据门控，而不是只由日期触发
    recovery_gate = np.clip((recovery_evidence - 0.45) / 0.35, 0.0, 1.0)
    stress_brake = np.clip(1.0 - 0.75 * physical_shift, 0.0, 1.0)
    conditional_easing = 0.75 * recovery_gate * stress_brake
    if day_index >= 120:
        logits[STATE_INDEX["easing"]] += 0.26 * conditional_easing
        logits[STATE_INDEX["escalation"]] -= 0.24 * conditional_easing
    if day_index >= 150:
        logits[STATE_INDEX["easing"]] += 0.18 * conditional_easing
        logits[STATE_INDEX["neutral"]] -= 0.07 * conditional_easing
        logits[STATE_INDEX["escalation"]] -= 0.11 * conditional_easing

    exp_values = np.exp(logits - np.max(logits))
    return exp_values / exp_values.sum()
\end{paperlisting}

代码~\ref{lst:state-transition-path} 展示状态转移路径的价格递推。每条样本路径先抽取下一期市场状态，再根据状态选择对应中心情景，并叠加历史波动率约束下的扰动和状态切换跳跃，由此得到长期条件区间和尾部风险。

\begin{paperlisting}{lst:state-transition-path}{python}{状态转移路径价格递推}
def simulate_path(rng, sample_id, dates, center_price_lookup, start_price,
                  daily_sigma, transition_jump_sigma_scale, risk_constraints,
                  physical_pressure_lookup, recovery_evidence_lookup):
    state = "neutral"
    previous_price = start_price
    previous_noise = 0.0
    rows = []

    for _, base_row in dates.iterrows():
        day_index = int(base_row["day_index"])
        physical_pressure = physical_pressure_lookup.get(day_index, 0.0)
        recovery_evidence = recovery_evidence_lookup.get(day_index, 0.0)
        next_state = sample_next_state(rng, state, day_index, risk_constraints,
                                       physical_pressure_lookup,
                                       recovery_evidence_lookup)
        spec = STATES[STATE_INDEX[next_state]]
        target_price = center_price_lookup[spec.center_scenario][day_index]

        noise_scale = max(previous_price * daily_sigma
                          * spec.volatility_multiplier, 0.15)
        innovation = rng.normal(0.0, noise_scale)
        previous_noise = 0.30 * previous_noise + innovation
        transition_jump = 0.0
        if next_state != state:
            transition_jump = rng.normal(spec.transition_jump_mean,
                                         spec.transition_jump_sigma
                                         * transition_jump_sigma_scale)

        price = 0.82 * previous_price + 0.18 * target_price
        price = np.clip(price + previous_noise + transition_jump, 70.0, 145.0)
        rows.append({"sample_id": sample_id, "day_index": day_index,
                     "trade_date": base_row["trade_date"], "状态": spec.label,
                     "forecast_price": price})
        state = next_state
        previous_price = price

    return pd.DataFrame(rows)
\end{paperlisting}

\clearpage
